% Pandoc 中文论文 LaTeX 模板
% 使用 ctexart 文档类，适配中文期刊排版
\documentclass[12pt,a4paper,UTF8]{ctexart}

% === 页面设置 ===
\usepackage[top=2.5cm, bottom=2.5cm, left=3cm, right=2.5cm]{geometry}
\usepackage{setspace}
\onehalfspacing

% === 字体（Windows 中文字体） ===
\setCJKmainfont{SimSun}[AutoFakeBold=2.5]
\setCJKsansfont{SimHei}
\setCJKmonofont{SimHei}

% === 表格 ===
\usepackage{booktabs}
\usepackage{longtable}
\usepackage{array}
\renewcommand{\arraystretch}{1.2}

% === 图表 ===
\usepackage{graphicx}
\usepackage{caption}
\usepackage{subcaption}
\graphicspath{{figures/}}

% === 参考文献 ===
\usepackage[backend=biber,style=gb7714-2015]{biblatex}
\addbibresource{references.bib}

% === 超链接 ===
\usepackage[hidelinks]{hyperref}

% === 数学 ===
\usepackage{amsmath,amssymb}

% === 页眉页脚 ===
\usepackage{fancyhdr}
\pagestyle{fancy}
\fancyhf{}
\fancyhead[C]{\small 《地书》图形语言的语言性量化评估}
\fancyfoot[C]{\thepage}
\renewcommand{\headrulewidth}{0.4pt}

% === 标题页 ===
\title{\textbf{《地书》图形语言的语言性量化评估\\——基于多维标注体系的统计分析}}
\author{池长俊\thanks{人工智能与计算机学院，江南大学} \quad 王顺祺 \quad 刘棪 \quad 高子阳}
\date{}

% === 摘要环境 ===
\newenvironment{chineseabstract}{%
  \begin{center}
    \begin{minipage}{0.9\textwidth}
      \small
      \textbf{摘\quad 要：}
}{%
    \end{minipage}
  \end{center}
  \vspace{0.5em}
}

% === 关键词 ===
\newcommand{\keywords}[1]{%
  \begin{center}
    \begin{minipage}{0.9\textwidth}
      \small
      \textbf{关键词：}#1
    \end{minipage}
  \end{center}
}

\begin{document}

\maketitle

\begin{chineseabstract}
图形符号系统（如Emoji、交通标志、UI图标）在当代数字通信中日益重要，但尚未有系统的计算框架来评估它们在多大程度上"像一种语言"。本文以徐冰《地书》为案例，提出一个多维度的"语言性"量化评估框架，涵盖频率结构、信息密度、词汇多样性、语法类别多样性、篇章结构复杂度和语义稳定性六个维度。基于708个标注文件中的46,494条语言学标注数据，对《地书》与八种自然语言语料库进行了系统对比。研究发现：（1）《地书》的词汇频率严格遵循齐夫幂律分布（R² = 0.923），但其Zipf斜率（−0.612）显著浅于中文自然语言（−1.35~−1.47），表明词频分布异常平坦；（2）在同等语料规模下，其词汇丰富度（降采样TTR = 0.434）约为中文的3倍，归一化熵（0.914）在所有对比语料中最高；（3）其语法类别分布呈现标点/边界标记主导（33.3\%）、动词性超过名词性（1.63:1）、虚词/语法标记极度稀缺（0.9\%）的三重特征；（4）篇章结构以时间先后为主线（43.7\%），复杂逻辑关系不足11\%。综合来看，《地书》在统计结构上遵循人类语言的深层组织规律，但其词汇分布模式反映了一个缺乏语法化功能词体系的符号系统所面临的系统性约束。本文首次为图形符号系统提供了一套可复用的多维度语言性评估方法，并将标注者间分歧重新定位为揭示符号系统固有歧义性的研究窗口。
\end{chineseabstract}

\keywords{地书；图形语言；语言性评估；标注者间一致性；齐夫定律；计算语言学}

\newpage

$body$

\end{document}
